% =====================================================================
%  Выводы по лабораторной работе.
% =====================================================================

\labpartbreak{Выводы}

В ходе выполнения лабораторной работы существующее приложение
\texttt{SecureApp} было расширено разделом, работающим исключительно с
оперативной памятью процесса, без какого-либо постоянного хранения, и
поддерживающим создание, обновление и удаление как конфиденциальных
(шифруемых непосредственно при вводе), так и неконфиденциальных данных.
Выполнен анализ кода на предмет обработки некорректного ввода (семь
проверок, ни одна не привела к падению или зависанию процесса), на
предмет потокобезопасности и на предмет обращения с конфиденциальными
данными в памяти (добавлена явная, но осознанно ограниченная -- в силу
неизменяемости строк в .NET -- очистка собственных промежуточных
буферов). Проведён анализ потребления процессорного времени и памяти под
нагрузкой средствами \texttt{dotnet-counters}, а также снято и
проанализировано три дампа оперативной памяти -- после создания, после
обновления и после удаления данных, -- с прямым поиском в них как
конфиденциальных, так и неконфиденциальных контрольных значений.

Семантически работа продемонстрировала принципиальную границу между тем,
что может гарантировать код приложения, и тем, что определяется средой
выполнения и фреймворком: приложение корректно не хранит конфиденциальные
данные в открытом виде ни в одной из своих собственных структур --
это подтверждено на уровне кода, -- однако открытый текст тем не менее
был обнаружен в дампе процесса на всех трёх контрольных точках, включая
дамп, снятый после удаления записи и принудительной полной сборки
мусора. Это не опровергает корректность реализованного шифрования, а
показывает границы его применимости: шифрование защищает данные там, где
код приложения ими управляет (в частности -- при попадании на диск, что
было отдельно проверено в лабораторной работе №2.1), но не может
защитить от копий, создаваемых инфраструктурой (буферизация HTTP-запроса
и ответа) ещё до и после того, как код приложения получает управление.
Разработчик, полагающийся исключительно на факт вызова функции
шифрования, рискует сделать ложный вывод о защищённости данных в памяти
процесса -- вывод, который не подтверждается прямым наблюдением
дампа.

С прагматической точки зрения работа показала, что стандартные
инструменты диагностики .NET (\texttt{dotnet-dump},
\texttt{dotnet-gcdump}, \texttt{dotnet-counters}) -- точные функциональные
аналоги перечисленных в задании средств JDK (\texttt{jvisualvm},
\texttt{jstack}, \texttt{jhat}) -- достаточны для содержательного
криминалистического анализа состояния оперативной памяти реального
веб-приложения без установки стороннего специализированного ПО, а
методика поиска контрольного маркера в необработанном дампе утилитой
\texttt{strings} (аналогичная применённой в лабораторной работе №2.1 к
файлу базы данных) одинаково хорошо работает и для файла, и для дампа
процесса. Также показано, что для защиты конфиденциальных данных в
памяти управляемой среды нельзя полагаться на класс-обёртку вроде
\texttt{SecureString} -- обоснование его неприменимости, а не просто его
отсутствие, само по себе является частью анализа безопасности кода.
